\documentclass[11pt,a4paper]{article}
\usepackage[utf8]{inputenc}
\usepackage[T1]{fontenc}
\usepackage[spanish]{babel}
\usepackage[margin=2.2cm]{geometry}
\usepackage{array,longtable,enumitem}
\usepackage{amsmath,amssymb}
\usepackage[colorlinks=true,urlcolor=blue,linkcolor=blue]{hyperref}
\setlength{\parindent}{0pt}
\setlength{\parskip}{5pt}
\setlength{\emergencystretch}{2em}
\newcommand{\escena}[4]{\textbf{Escena #1: #2.} De #3 a #4}
\hypersetup{pdftitle={Guion de video EID MAT1188, Grupo 6},pdfauthor={Grupo 6}}
\begin{document}
\begin{center}
{\Large\textbf{Guion de video: EID Cálculo Intermedio}}\\[4pt]
{\large MAT1188. Grupo 6}\\[4pt]
Integral definida y área geométrica: contribuciones positivas y negativas\\[5pt]
\textbf{Duración propuesta: 5 minutos y 15 segundos.}
\end{center}

\section*{Cómo usar este guion}
El texto de narración está escrito para leerlo en voz alta. Las fórmulas y las indicaciones de pantalla sirven como apoyo y no se leen de principio a fin. Los tiempos son una propuesta que se debe comprobar mediante ensayo.

El montaje revisado anteriormente dura 2:56 y no tiene audio. Para preparar la versión final hay que incorporar las tres voces y ajustar las escenas a lo que se explica. Cada intervención debe quedar identificada. El video completo, incluidos los créditos, debe durar entre 4 y 6 minutos.

Los archivos del proyecto están en \texttt{scenes/escenas.py} y \texttt{utils/comun.py}. Los apoyos indicados como pendientes deben incorporarse y revisarse antes de grabar la versión definitiva.

\section*{Datos y reparto}
\begin{itemize}[leftmargin=1.5em,itemsep=3pt]
\item \textbf{Pregunta central:} ¿Por qué una integral definida puede ser menor que el área geométrica total?
\item \textbf{Ejemplo:} $f(x)=x^2-4$ en $[-3,3]$, con raíces $-2$ y $2$.
\item \textbf{Resultados:} $A^+=14/3$, $A^-=32/3$, integral $-6$ y área total $46/3$.
\item \textbf{Catalina Javiera Vergara Alarcón:} fundamentos, de 00:00 a 01:45.
\item \textbf{Ignacio Alexander De Celis Valencia:} ejemplo y aproximación, de 01:45 a 03:30.
\item \textbf{Marcelo Fabián Santana Astorga:} límite, resultados y cierre, de 03:30 a 05:15.
\end{itemize}

Antes de grabar, comprobar que el ejemplo sea distinto al del apunte base. Los tres deben comprender el trabajo completo para la defensa.

\clearpage
\section*{Catalina: bloque 1}

\noindent\begin{minipage}[t]{\textwidth}
\subsection*{\escena{1}{Portada}{00:00}{00:15}}

\textbf{Tiempo previsto:} 15 segundos. \textbf{Narración:} 23 palabras.

\subsubsection*{Qué se ve en pantalla}
Curso, grupo, integrantes y pregunta central. El rótulo del bloque debe identificar a Catalina.

\subsubsection*{Narración}
Hola, somos Catalina, Ignacio y Marcelo, del grupo seis. Explicaremos por qué una integral definida puede ser menor que el área geométrica total.

\subsubsection*{Indicaciones para el montaje}
Mostrar la presentación y dejar visible la pregunta al terminar. No es necesario leer todos los datos de la portada.
\end{minipage}\par\medskip

\noindent\begin{minipage}[t]{\textwidth}
\subsection*{\escena{2}{TeoriaSigno}{00:15}{01:00}}

\textbf{Tiempo previsto:} 45 segundos. \textbf{Narración:} 95 palabras.

\subsubsection*{Qué se ve en pantalla}
Rectángulos sobre y bajo el eje. Después, mostrar la suma y sus elementos:
\[
 S_n=\sum_{i=1}^{n}f(x_i^*)\Delta x,\qquad
 \Delta x=\frac{b-a}{n},\qquad a<b.
\]
El dibujo inicial es ilustrativo. El ejemplo principal se presenta en la escena 4.

\subsubsection*{Narración}
Soy Catalina. Una suma de Riemann nos permite sumar pequeños aportes de una función. Dividimos el intervalo en n partes del mismo ancho, llamado delta x. En cada parte elegimos un punto muestra, que indica dónde calculamos la altura del rectángulo.

El aporte es esa altura multiplicada por el ancho. Como el ancho es positivo, el signo depende de la altura. Sobre el eje, el aporte es positivo. Bajo el eje, es negativo. Para una función continua en el intervalo, estas sumas se acercan a la integral definida al usar rectángulos cada vez más angostos.

\subsubsection*{Indicaciones para el montaje}
Señalar el ancho, el punto muestra y la altura al nombrarlos. Mostrar luego un aporte positivo y uno negativo. Dejar una pausa para leer la suma.
\end{minipage}\par\medskip

\noindent\begin{minipage}[t]{\textwidth}
\subsection*{\escena{3}{TeoriaDefiniciones}{01:00}{01:45}}

\textbf{Tiempo previsto:} 45 segundos. \textbf{Narración:} 96 palabras.

\subsubsection*{Qué se ve en pantalla}
Mostrar las dos áreas, ambas no negativas, y después las operaciones:
\[
 I=\int_a^b f(x)\,dx=A^+-A^-,\qquad
 A_{\mathrm{total}}=\int_a^b |f(x)|\,dx=A^++A^-.
\]
La relación con el valor absoluto es un apoyo que se debe incorporar si no aparece todavía en las escenas.

\subsubsection*{Narración}
Llamamos A más al área de las regiones sobre el eje y A menos al área de las regiones bajo él. Las dos son cantidades positivas o cero. El nombre A menos no significa que esa área sea negativa.

Para obtener la integral, restamos A menos a A más. Para obtener el área geométrica total, sumamos ambas. Esto también se expresa integrando el valor absoluto de la función, porque así contamos todos los aportes como positivos. La diferencia está en que la integral permite que los aportes se compensen. Ignacio mostrará cómo ocurre en nuestro ejemplo.

\subsubsection*{Indicaciones para el montaje}
Relacionar cada región con su símbolo antes de mostrar las operaciones. Leer las fórmulas con calma y pasar al ejemplo.
\end{minipage}\par\medskip

\clearpage
\section*{Ignacio: bloque 2}

\noindent\begin{minipage}[t]{\textwidth}
\subsection*{\escena{4}{EjemploPrincipal}{01:45}{02:15}}

\textbf{Tiempo previsto:} 30 segundos. \textbf{Narración:} 63 palabras.

\subsubsection*{Qué se ve en pantalla}
Identificar a Ignacio. Mostrar $f(x)=x^2-4$ en $[-3,3]$ y las raíces $x=-2$ y $x=2$. La función es positiva en $[-3,-2)$ y $(2,3]$, negativa en $(-2,2)$ y cero en las raíces.

\subsubsection*{Narración}
Soy Ignacio. Elegimos la función x al cuadrado menos cuatro, en el intervalo de menos tres a tres. Es un polinomio, así que es continua. Al igualarla a cero encontramos las raíces menos dos y dos.

La función es positiva en los tramos exteriores y negativa entre las raíces. En las raíces vale cero. Por eso distinguimos dos regiones verdes y una roja.

\subsubsection*{Indicaciones para el montaje}
Mostrar primero función e intervalo. Señalar las raíces y luego colorear las regiones. Confirmar antes de grabar que el ejemplo sea distinto al del apunte.
\end{minipage}\par\medskip

\noindent\begin{minipage}[t]{\textwidth}
\subsection*{\escena{5}{Rectangulos}{02:15}{02:50}}

\textbf{Tiempo previsto:} 35 segundos. \textbf{Narración:} 79 palabras.

\subsubsection*{Qué se ve en pantalla}
Partición con $n=6$, ancho $\Delta x=1$ y puntos medios. Puntos muestra: $-2{,}5$, $-1{,}5$, $-0{,}5$, $0{,}5$, $1{,}5$ y $2{,}5$. Mostrar altura y aporte.

\subsubsection*{Narración}
Dividimos el intervalo en seis partes iguales y elegimos el punto medio de cada una. Como el intervalo mide seis unidades, el ancho de cada rectángulo es uno.

La altura se obtiene evaluando la función en el punto medio. Por ejemplo, en dos coma cinco, la altura es dos coma veinticinco. Al multiplicarla por el ancho, obtenemos un aporte de dos coma veinticinco. Usamos el mismo procedimiento para los otros cinco rectángulos y conservamos el signo de cada aporte.

\subsubsection*{Indicaciones para el montaje}
Presentar la partición, los puntos medios y los rectángulos en ese orden. Hacer coincidir el ejemplo numérico con el rectángulo indicado.
\end{minipage}\par\medskip

\noindent\begin{minipage}[t]{\textwidth}
\subsection*{\escena{6}{SumaDeAportes}{02:50}{03:30}}

\textbf{Tiempo previsto:} 40 segundos. \textbf{Narración:} 91 palabras.

\subsubsection*{Qué se ve en pantalla}
Mostrar por etapas:
\[
 S_6=4{,}5-11=-6{,}5,\qquad
 A_6=4{,}5+11=15{,}5.
\]
Las dos cantidades son aproximaciones. Los cuatro rectángulos negativos pertenecen a una misma región bajo la curva.

\subsubsection*{Narración}
Los dos rectángulos positivos aportan dos coma veinticinco cada uno. En total suman cuatro coma cinco. Los cuatro negativos aportan menos uno coma setenta y cinco, menos tres coma setenta y cinco, menos tres coma setenta y cinco y menos uno coma setenta y cinco. Su suma es menos once.

Por eso, la suma con signo da menos seis coma cinco. Para aproximar el área, sumamos los valores sin signo y obtenemos quince coma cinco. Son dos aproximaciones hechas con los mismos rectángulos. Marcelo explicará qué pasa al aumentar su número.

\subsubsection*{Indicaciones para el montaje}
Mostrar cada suma mientras se explica. Mantener visible la palabra aproximación y dejar una pausa antes del cambio de voz.
\end{minipage}\par\medskip

\clearpage
\section*{Marcelo: bloque 3}

\noindent\begin{minipage}[t]{\textwidth}
\subsection*{\escena{7}{Refinamiento}{03:30}{03:55}}

\textbf{Tiempo previsto:} 25 segundos. \textbf{Narración:} 63 palabras.

\subsubsection*{Qué se ve en pantalla}
Identificar a Marcelo. Pasar de $n=6$ a $n=12$ y $n=24$. La tabla muestra la suma con signo y el área aproximada.

\subsubsection*{Narración}
Soy Marcelo. Conservamos la función, el intervalo y la elección del punto medio, pero aumentamos el número de rectángulos. Su ancho disminuye y la aproximación mejora en este ejemplo.

La suma con signo se acerca a menos seis. La suma de los valores sin signo se acerca a quince coma treinta y tres. Este proceso nos permite entender la integral como un límite.

\subsubsection*{Indicaciones para el montaje}
Hacer coincidir el cambio de n con la voz. No leer todos los números de la tabla. Si es necesario aclararlo, n tiende a infinito como proceso, no como un valor que se sustituye.
\end{minipage}\par\medskip

\noindent\begin{minipage}[t]{\textwidth}
\subsection*{\escena{8}{ComparacionFinal}{03:55}{04:45}}

\textbf{Tiempo previsto:} 50 segundos. \textbf{Narración:} 95 palabras.

\subsubsection*{Qué se ve en pantalla}
Añadir el cálculo exacto antes de las barras y flechas. Mostrar una línea por vez:
\[
 F(x)=\frac{x^3}{3}-4x,\qquad F'(x)=x^2-4.
\]
\[
 A^+=[F(x)]_{-3}^{-2}+[F(x)]_{2}^{3}
 =\frac{7}{3}+\frac{7}{3}=\frac{14}{3}.
\]
\[
 \int_{-2}^{2}f(x)\,dx=-\frac{32}{3},\qquad A^-=\frac{32}{3}.
\]
\[
 I=\frac{14}{3}-\frac{32}{3}=-6,\qquad
 A_{\mathrm{total}}=\frac{14}{3}+\frac{32}{3}=\frac{46}{3}.
\]
Este desarrollo es un apoyo previsto en el guion. Hay que comprobar su incorporación al video.

\subsubsection*{Narración}
Para comprobar los valores exactos usamos una antiderivada. Es una función cuya derivada devuelve la función original. En este caso es x al cubo dividido por tres, menos cuatro x. La evaluamos en los extremos de cada tramo y restamos los resultados.

Cada región exterior tiene área siete tercios. Juntas dan catorce tercios. En la región central, la integral vale menos treinta y dos tercios y su área es treinta y dos tercios. Al restar las áreas obtenemos la integral, menos seis. Al sumarlas obtenemos cuarenta y seis tercios, aproximadamente quince coma treinta y tres.

\subsubsection*{Indicaciones para el montaje}
Explicar primero el cálculo exacto. Luego señalar cómo las barras suman ambas áreas y cómo las flechas avanzan y retroceden hasta llegar a menos seis. Dejar pausas para seguir las operaciones.
\end{minipage}\par\medskip

\noindent\begin{minipage}[t]{\textwidth}
\subsection*{\escena{9}{RespuestaCentral}{04:45}{05:05}}

\textbf{Tiempo previsto:} 20 segundos. \textbf{Narración:} 40 palabras.

\subsubsection*{Qué se ve en pantalla}
Pregunta central y comparación final. Mostrar la condición $A^->0$ para la desigualdad estricta.

\subsubsection*{Narración}
Las regiones bajo el eje aportan valores negativos a la integral. El área geométrica suma todas las áreas sin signo. Por eso, si A menos es mayor que cero, la integral es menor. Si A menos vale cero, ambas coinciden.

\subsubsection*{Indicaciones para el montaje}
Responder mirando la comparación, sin volver a leer todas las fórmulas. Aclarar que menos seis no representa un área geométrica negativa.
\end{minipage}\par\medskip

\noindent\begin{minipage}[t]{\textwidth}
\subsection*{\escena{10}{Creditos}{05:05}{05:15}}

\textbf{Tiempo previsto:} 10 segundos. \textbf{Narración:} 25 palabras.

\subsubsection*{Qué se ve en pantalla}
Mostrar el apunte, Stewart, Thomas y Manim. Incluir Claude y Codex, sus usos y la forma de comprobar los resultados. El detalle puede quedar en la ficha para que se lea sin apuro. Comprobar qué fuentes y recursos se utilizaron realmente.

\subsubsection*{Narración}
Las fuentes y el uso de inteligencia artificial están detallados en los créditos y en la ficha. Comprobamos los cálculos y revisamos las escenas. Gracias.

\subsubsection*{Indicaciones para el montaje}
Mantener los créditos legibles. Evitar comprimir toda la ficha en una sola pantalla o leer la bibliografía completa en voz alta.
\end{minipage}\par\medskip

\clearpage
\section*{Resumen de tiempos}
\begin{longtable}{|p{0.8cm}|p{4.0cm}|c|c|c|p{1.8cm}|}
\hline
\textbf{N.} & \textbf{Escena} & \textbf{Inicio} & \textbf{Fin} & \textbf{Seg.} & \textbf{Voz}\\\hline

1 & Portada & 00:00 & 00:15 & 15 & Catalina \\\hline

2 & TeoriaSigno & 00:15 & 01:00 & 45 & Catalina \\\hline

3 & TeoriaDefiniciones & 01:00 & 01:45 & 45 & Catalina \\\hline

4 & EjemploPrincipal & 01:45 & 02:15 & 30 & Ignacio \\\hline

5 & Rectangulos & 02:15 & 02:50 & 35 & Ignacio \\\hline

6 & SumaDeAportes & 02:50 & 03:30 & 40 & Ignacio \\\hline

7 & Refinamiento & 03:30 & 03:55 & 25 & Marcelo \\\hline

8 & ComparacionFinal & 03:55 & 04:45 & 50 & Marcelo \\\hline

9 & RespuestaCentral & 04:45 & 05:05 & 20 & Marcelo \\\hline

10 & Creditos & 05:05 & 05:15 & 10 & Marcelo \\\hline

\multicolumn{4}{|r|}{\textbf{Total propuesto}} & \textbf{315} & \textbf{5:15}\\\hline
\end{longtable}

\section*{Revisión antes de grabar}
\begin{itemize}[leftmargin=1.5em,itemsep=4pt]
\item Ensayar las tres voces y comprobar los tiempos reales. No hacer un corte si la explicación todavía no termina.
\item Dar más tiempo al contenido matemático que a la presentación y los créditos.
\item Añadir los apoyos del valor absoluto y del cálculo exacto si todavía no están en las escenas. Revisar que se lean bien.
\item Mantener la diferencia entre aproximación y valor exacto. Con seis rectángulos, $S_6=-6{,}5$. La integral exacta vale $-6$.
\item Comprobar el apunte y las dos fuentes matemáticas adicionales. No añadir números de página sin revisar la edición consultada.
\item Actualizar los créditos y la ficha con el uso real de Claude y Codex. Incluir también el apoyo en el guion y en la revisión de la ficha.
\item Solo declarar grabación y revisión final del audio después de realizarlas.
\item Completar sección, duración real, enlace al video y versión del código en la ficha de entrega.
\end{itemize}

\section*{Preparación para la defensa}
La rúbrica indica un máximo de diez minutos y la elección al azar de un integrante presente. La persona elegida responde cuatro preguntas. Por eso, cada integrante debe poder explicar el tema, justificar el procedimiento, interpretar los resultados y describir cómo afectan los parámetros del código al dibujo. No hace falta memorizar instrucciones de Manim.

\end{document}
